\section{USAMO 1982}
        \begin{enumerate}
        	\item (\textit{USAMO 1982}) In a party with $1982$ persons, among any group of four there is at least one person who knows each of the other three. What is the minimum number of people in the party who know everyone else?


 

	\item (\textit{USAMO 1982}) Let $S_r=x^r+y^r+z^r$ with $x,y,z$ real. It is known that if $S_1=0$,


 $(*)$ $\frac{S_{m+n}}{m+n}=\frac{S_m}{m}\frac{S_n}{n}$


 for $(m,n)=(2,3),(3,2),(2,5)$, or $(5,2)$. Determine \textit{all} other pairs of integers $(m,n)$ if any, so that $(*)$ holds for all real numbers $x,y,z$ such that $S_1=0$.


 

	\item (\textit{USAMO 1982}) If a point $A_1$ is in the interior of an equilateral triangle $ABC$ and point $A_2$ is in the interior of $\triangle{A_1BC}$, prove that


 $I.Q. (A_1BC) > I.Q.(A_2BC)$,


 where the \textit{isoperimetric quotient} of a figure $F$ is defined by


 $I.Q.(F) = \frac{\text{Area (F)}}{\text{[Perimeter (F)]}^2}$


 

	\item (\textit{USAMO 1982}) Prove that there exists a positive integer $k$ such that $k\cdot2^n+1$ is composite for every positive integer $n$.


 

	\item (\textit{USAMO 1982}) $A,B$, and $C$ are three interior points of a sphere $S$ such that $AB$ and $AC$ are perpendicular to the diameter of $S$ through $A$, and so that two spheres can be constructed through $A$, $B$, and $C$ which are both tangent to $S$. Prove that the sum of their radii is equal to the radius of $S$.


 

\end{enumerate}